\documentclass[10pt,twocolumn]{article}

\usepackage[margin=0.75in]{geometry}
\usepackage[utf8]{inputenc}
\usepackage[T1]{fontenc}
\usepackage{booktabs}
\usepackage{graphicx}
\usepackage{amsmath}
\usepackage{url}
\usepackage[hidelinks]{hyperref}
\usepackage{siunitx}
\usepackage{caption}
\usepackage{balance}

\sisetup{group-separator={,}}
\captionsetup{font=small,labelfont=bf}

% Table bodies are generated files read from inside a tabular. LaTeX's \input
% injects a \par at end of file, which breaks the alignment before \bottomrule;
% the TeX primitive does not.
\makeatletter
\newcommand{\inputrows}[1]{\@@input #1 }
\makeatother

% Numbers regenerated from research/results/*.json by make_macros.py.
\IfFileExists{generated/results.tex}{% Generated by research/experiments/make_macros.py -- do not edit.
\newcommand{\ResTrials}{100}
\newcommand{\ResWarmupDays}{7}
\newcommand{\ResMeasureDays}{7}
\newcommand{\ResEventsPerWindow}{5,674}
\newcommand{\ResFPDay}{0.45}
\newcommand{\ResFPDayCI}{0.04}
\newcommand{\ResPrecision}{83.5}
\newcommand{\ResTTDAone}{47.6}
\newcommand{\ResTTCAone}{67.6}
\newcommand{\ResDetAone}{100}
\newcommand{\ResTTDMedianAone}{48.0}
\newcommand{\ResTTDPNinetyFiveAone}{48.0}
\newcommand{\ResTTDMaxAone}{48.0}
\newcommand{\ResCIAone}{0.5}
\newcommand{\ResFirstRuleAone}{R01}
\newcommand{\ResFirstRuleTrialsAone}{100}
\newcommand{\ResFallbackTrialsAone}{0}
\newcommand{\ResFallbackRuleAone}{--}
\newcommand{\ResTTDAtwo}{1785.9}
\newcommand{\ResTTCAtwo}{1824.3}
\newcommand{\ResDetAtwo}{100}
\newcommand{\ResTTDMedianAtwo}{4.0}
\newcommand{\ResTTDPNinetyFiveAtwo}{5924.0}
\newcommand{\ResTTDMaxAtwo}{49544.0}
\newcommand{\ResCIAtwo}{1673.6}
\newcommand{\ResFirstRuleAtwo}{R02}
\newcommand{\ResFirstRuleTrialsAtwo}{92}
\newcommand{\ResFallbackTrialsAtwo}{8}
\newcommand{\ResFallbackRuleAtwo}{R05}
\newcommand{\ResTTDAthree}{13.6}
\newcommand{\ResTTCAthree}{33.5}
\newcommand{\ResDetAthree}{100}
\newcommand{\ResTTDMedianAthree}{14.0}
\newcommand{\ResTTDPNinetyFiveAthree}{14.0}
\newcommand{\ResTTDMaxAthree}{14.0}
\newcommand{\ResCIAthree}{0.3}
\newcommand{\ResFirstRuleAthree}{R06}
\newcommand{\ResFirstRuleTrialsAthree}{100}
\newcommand{\ResFallbackTrialsAthree}{0}
\newcommand{\ResFallbackRuleAthree}{--}
\newcommand{\ResTTDAfour}{4.0}
\newcommand{\ResTTCAfour}{24.0}
\newcommand{\ResDetAfour}{100}
\newcommand{\ResTTDMedianAfour}{4.0}
\newcommand{\ResTTDPNinetyFiveAfour}{4.0}
\newcommand{\ResTTDMaxAfour}{4.0}
\newcommand{\ResCIAfour}{0.0}
\newcommand{\ResFirstRuleAfour}{R08}
\newcommand{\ResFirstRuleTrialsAfour}{93}
\newcommand{\ResFallbackTrialsAfour}{7}
\newcommand{\ResFallbackRuleAfour}{R05}
\newcommand{\ResTTDAfive}{242.1}
\newcommand{\ResTTCAfive}{262.8}
\newcommand{\ResDetAfive}{100}
\newcommand{\ResTTDMedianAfive}{246.0}
\newcommand{\ResTTDPNinetyFiveAfive}{246.0}
\newcommand{\ResTTDMaxAfive}{246.0}
\newcommand{\ResCIAfive}{2.0}
\newcommand{\ResFirstRuleAfive}{R09}
\newcommand{\ResFirstRuleTrialsAfive}{99}
\newcommand{\ResFallbackTrialsAfive}{1}
\newcommand{\ResFallbackRuleAfive}{R10}
\newcommand{\ResTTDMean}{418.6}
\newcommand{\ResTTDStream}{4.0}
\newcommand{\AblFullCrit}{0.8}
\newcommand{\AblNoSDLSCrit}{2.8}
\newcommand{\AblNoReplayCrit}{2.0}
\newcommand{\AblNoneCrit}{6.6}
\newcommand{\AblNoneMB}{13,086}
\newcommand{\AblNoResponseMB}{13,086}
\newcommand{\AblFullMB}{1,213}
\newcommand{\AblFullSucc}{51.7}
\newcommand{\AblNoResponseSucc}{288.0}
\newcommand{\SweepIntegOneFP}{1.13}
\newcommand{\SweepIntegOneTTD}{60.5}
\newcommand{\SweepIntegThreeFP}{0.50}
\newcommand{\SweepIntegThreeTTD}{62.5}
\newcommand{\SweepBatchFastTTD}{38.5}
\newcommand{\SweepBatchSlowTTD}{212.5}
\newcommand{\CostMonthly}{419}
\newcommand{\CostTopService}{OpenSearch Serverless}
\newcommand{\CostTopShare}{82}
\newcommand{\CostPerHundred}{79.89}
}{%
  \typeout{WARNING: run `make macros` first; using placeholders}%
  \newcommand{\ResTrials}{??}\newcommand{\ResWarmupDays}{??}%
  \newcommand{\ResMeasureDays}{??}\newcommand{\ResEventsPerWindow}{??}%
  \newcommand{\ResFPDay}{??}\newcommand{\ResFPDayCI}{??}%
  \newcommand{\ResPrecision}{??}\newcommand{\ResTTDMean}{??}%
  \newcommand{\ResTTDStream}{??}\newcommand{\CostMonthly}{??}%
}

\title{\textbf{A Cloud-Native Security Architecture for Detecting Cyber Threats
in Satellite Ground Segments}}

\author{%
  \normalsize Space Ground Station Security Lab\\
  \normalsize \url{https://github.com/rikosoo/space-ground-station-security-lab}
}

\date{}

\begin{document}
\maketitle

\begin{abstract}
Satellite ground segments have become ordinary cloud workloads while retaining
an extraordinary property: their actuator is a vehicle that cannot be rebooted,
re-imaged, or taken offline. Security guidance for this domain is plentiful and
largely qualitative; measurements of whether cloud-native security telemetry
actually detects attacks against a ground segment are not. We ask a single
question --- \emph{can cloud-native security telemetry effectively detect
cyberattacks against satellite ground segments?} --- and answer it with a
reproducible testbed rather than an architecture diagram. The testbed simulates
the full chain from spacecraft to incident response, executes five controlled
attacks (credential compromise, unauthorized telecommand, telemetry tampering,
telecommand replay, and mission-data exfiltration), and measures detection
latency, false positives, defensive-control attribution, and cost. Across
\ResTrials{} trials all five attacks are detected, with a mean time to detect of
\ResTTDMean{}\,s and \ResFPDay{}~$\pm$~\ResFPDayCI{} false positives per day at
\ResPrecision{}\,\% precision, for a modelled cost of approximately
USD~\CostMonthly{} per month for a single spacecraft. Three findings generalise
beyond the testbed: detection latency is set by the tier a rule runs in rather
than by the attack; cryptographic frame integrity and physics-aware analytics are
complementary rather than redundant, because an adversary who disables the former
is only visible to the latter; and reporting detection without reporting what the
response failed to prevent overstates the security of the system. All code,
detection content, labelled data and analysis scripts are released.
\end{abstract}

\section{Introduction}

A modern satellite ground segment is, from the outside, a fairly ordinary cloud
workload: an authenticated web API, a message stream, an object store, and a
handful of serverless functions. From the inside it is not ordinary at all. Its
actuator is a spacecraft that cannot be rebooted, re-imaged, physically
inspected, or taken offline for maintenance. It is reachable for a few minutes a
day --- in the configuration studied here, three passes totalling roughly
\SI{22}{\minute} out of every \SI{24}{\hour}, leaving the adversary the other
\SI{97}{\percent} of the day to work in. And the consequence of a successful
attack is not data loss but the loss of a physical asset in orbit.

The security guidance for this domain has grown quickly and is largely
qualitative. Policy sets principles \cite{spd5}; introductory guidance describes
the problem space \cite{nistir8270}; threat knowledge bases enumerate techniques
\cite{sparta}; the research literature has established that the attack surface is
real, both by analysis \cite{falco2019principles,manulis2021newspace,pavur2020launchpad}
and by finding exploitable flaws in flight software actually used in orbit
\cite{willbold2023odyssey}. The 2022 KA-SAT incident \cite{viasat2022}
demonstrated that ground-segment compromise, rather than exotic attacks on the
space segment, is the practical path to mission impact.

What is scarce is measurement. Given a ground segment instrumented the way a
competent team would instrument any cloud workload --- an audit trail, a
streaming detection tier, findings aggregation, automated containment --- how
quickly are attacks actually detected? Which log makes each one visible? How much
noise do the same detections produce during a week in which nothing happens? Which
defensive mechanism did the work: the one that produced the alert, or one that
quietly prevented the attack from mattering? And what does the whole arrangement
cost to run?

\subsection*{Research question and sub-questions}

\begin{quote}
\textbf{RQ.} Can cloud-native security telemetry effectively detect cyberattacks
against satellite ground segments?
\end{quote}

We decompose it into five measurable sub-questions:

\begin{description}
\item[RQ1 (Detection).] Are the five attack classes detected, and how long does
detection take relative to the first malicious action?
\item[RQ2 (Noise).] How many false positives do the same detections produce
during benign operations, and at what precision?
\item[RQ3 (Evidence).] Which log source and which rule identify each attack?
\item[RQ4 (Attribution).] Which defensive mechanism prevents or limits each
attack, as opposed to merely observing it?
\item[RQ5 (Cost).] What does the architecture cost per month, and how does that
cost scale with fleet size?
\end{description}

\subsection*{Contributions}

\begin{enumerate}
\item \textbf{A reproducible testbed} spanning spacecraft dynamics, RF link,
ground station, mission-control API, cloud analytics and automated response,
running deterministically on a virtual clock: fourteen simulated days in about
two seconds, on a laptop, with no cloud account.
\item \textbf{Five controlled attack scenarios} with explicit adversary
capability levels and per-event ground-truth labels, enabling true positives to
be defined causally rather than by temporal coincidence.
\item \textbf{A detection pack of eleven rules} in two synchronised forms:
executable Python that both the experiments and the deployed Lambda import, and
portable Sigma \cite{sigma} for other backends.
\item \textbf{An empirical evaluation} of detection latency, false positives,
per-rule precision, a defence ablation that attributes outcomes to individual
controls, a sensitivity analysis of the tuning parameters, and a cost model.
\item \textbf{A labelled dataset} and the full analysis pipeline, so every number
in this paper can be regenerated with one command.
\end{enumerate}

\section{Background and related work}

\subsection{The space link and its security services}

Telecommand and telemetry in most civil and commercial missions follow the CCSDS
protocol stack: application data in Space Packets \cite{ccsds133spp}, carried in
transfer frames, optionally protected by the Space Data Link Security protocol
(SDLS) \cite{ccsds355sdls}. SDLS provides authentication, optional confidentiality,
and --- critically for this work --- an anti-replay sequence number bound into the
authenticated data. Our testbed models SDLS at the level of abstraction the
detections care about: a keyed authentication tag over the frame payload and
sequence number, and a sliding acceptance window on the vehicle. It does not
model bit layouts, which do not change which events reach the analytics tier.

The two properties that matter for detection engineering are worth stating
plainly. First, a frame that fails its tag check on the vehicle is an
unambiguous signal: nominal operations never produce one, because the ground
system signs every frame with the key the vehicle holds. Second, authentication
alone does not establish freshness --- a recorded frame remains cryptographically
valid forever --- which is why the anti-replay window, not the tag, is the
control that defeats replay.

\subsection{Threat knowledge for space systems}

SPARTA \cite{sparta} adapts the ATT\&CK \cite{attack} structure to space systems,
covering reconnaissance through impact for both the space and ground segments. We
map every detection rule to a SPARTA technique and, where an analogue exists, to
an ATT\&CK technique; this makes the coverage claim checkable rather than
rhetorical. NIST IR 8270 \cite{nistir8270} and SPD-5 \cite{spd5} supply the
control expectations --- authentication of commands, protection of telemetry
integrity, monitoring --- that the architecture implements.

\subsection{Prior empirical work}

Pavur and Martinovic \cite{pavur2020launchpad} argue that satellite security
research is held back by the absence of accessible testbeds, and that the field
substitutes speculation for measurement as a result. Willbold et al.
\cite{willbold2023odyssey} answer part of that challenge on the space segment,
analysing real satellite firmware and finding that several vehicles accept
unauthenticated commands outright. Manulis et al. \cite{manulis2021newspace}
survey the New Space threat landscape and identify the ground segment as the
dominant practical attack surface, a judgement the KA-SAT incident
\cite{viasat2022} supports.

Our work is complementary and deliberately narrower. We do not analyse flight
software. We assume the space segment implements its security services correctly
and ask what the \emph{ground} segment's cloud telemetry sees when an adversary
attacks the chain around it. To our knowledge this combination --- an end-to-end
ground-segment testbed, five labelled attack scenarios, and reported detection
latency, false-positive rate, control attribution and cost --- has not previously
been published in a form a reader can rerun.

\subsection{Why detection latency is the right metric here}

In enterprise security, dwell time is measured in days and the marginal value of
shaving seconds is low. A ground segment inverts this. The adversary's window of
\emph{effect} is the pass: a command that cannot be transmitted cannot harm the
vehicle. A pass in our configuration lasts about \SI{7}{\minute}. A detection
that fires in seconds can contain an intrusion before the next pass; one that
fires in an hour cannot. Latency is therefore not a quality-of-service metric in
this domain --- it decides whether the response is capable of mattering.

\section{Threat model}

\subsection{System under consideration}

A single spacecraft in a \SI{550}{\kilo\metre}, \SI{53}{\degree}-inclination
orbit (period \SI{95.7}{\minute}, eclipse fraction \num{0.37}), operated through
one ground station at \SI{-23.55}{\degree} latitude with a
\SI{10}{\degree} mask angle. This geometry yields three passes per day averaging
\SI{437}{\second} each. Mission control, the telemetry archive and all security
analytics run in one cloud account.

\subsection{Assets}

\begin{itemize}
\item \textbf{Telecommand authority.} Direct control of a physical vehicle;
worst case is loss of the spacecraft.
\item \textbf{Frame authentication keys.} Possession allows indefinite forgery
that no tag check can detect.
\item \textbf{Telemetry integrity.} The danger is not the falsified number but
the operator decision taken because of it.
\item \textbf{The telemetry archive.} Payload products, orbit determination and
mission intelligence.
\item \textbf{The audit trail.} Cover for everything above.
\end{itemize}

\subsection{Adversary capability levels}

Each scenario declares which level it assumes, so that results are not
implicitly claimed against a stronger or weaker adversary than the one executed.

\begin{description}
\item[N1 --- remote, credentialed.] Phishing, credential stuffing, relay of a
second factor. No radio capability. Covers A1, A5 and the API portion of A2.
\item[N2 --- radio-capable.] Can receive and transmit on the mission's bands and
knows the frame format. Cannot break HMAC-SHA256. Covers A4 and the RF portion
of A2.
\item[N3 --- ground-segment foothold.] Controls a host between the demodulator
and cloud ingestion. Covers A3.
\end{description}

Out of scope by assumption: an adversary who obtains the SDLS key, breaks the
cryptography, or compromises the cloud provider's control plane. Each defeats the
architecture by definition and would make the measurement uninformative rather
than pessimistic. Also out of scope: jamming and other availability attacks,
which are detected by link metrics rather than security telemetry, and whose
inclusion would confound the false-positive measurement.

\subsection{Assumptions}

\begin{enumerate}
\item Frame authentication keys are generated in, and never leave, a managed key
store.
\item The vehicle enforces its anti-replay window; the ground segment cannot
persuade it to accept a stale frame.
\item Telemetry is archived before it is analysed, so detection never depends on
the analytics tier being available in real time.
\item The detection pipeline is itself monitored. A silent pipeline is treated as
an outage, not as a quiet week --- an alarm on consumer lag is part of the
architecture, not an operational nicety.
\item The operator population is small and enumerable. This is what makes
network-origin profiling substantially more effective in a ground segment than
in a general enterprise.
\end{enumerate}

\section{Architecture}

\subsection{Layers}

The testbed implements five layers, each emitting events into one normalized
schema:

\begin{enumerate}
\item \textbf{Spacecraft.} Circular-orbit propagation, power/thermal/attitude
dynamics with physically bounded rates, a thirteen-command catalogue with a
\emph{critical} flag, SDLS tag verification and an anti-replay window.
\item \textbf{Ground station.} Pass geometry, an elevation-dependent frame error
model, integrity verification of downlinked frames, and normalization of every
observation into the event schema.
\item \textbf{Mission-control API.} Authentication with optional second factor,
role-based command allowlists, per-principal rate limiting, and a pass-window
feasibility check. Frames are signed here before transmission.
\item \textbf{Cloud.} A stream carrying normalized events, an encrypted object
store for telemetry products, an audit trail covering both management and object
data events, a streaming detection function, findings aggregation, metrics and
alarms.
\item \textbf{Response.} A containment function subscribed to high and critical
findings, implementing the containment step --- and only the containment step ---
of the incident-response playbooks.
\end{enumerate}

\subsection{Two paths, two latencies}

The architecture has exactly two detection paths and their latencies differ by an
order of magnitude:

\begin{description}
\item[Stream path.] Ground-segment observations (authentication, telecommands,
frame integrity, telemetry samples) reach the detection function through the
stream within seconds. Its bound is the consumer's batching window.
\item[Control-plane path.] Cloud API calls reach the same function through the
audit trail and the event bus. Its bound is trail delivery, which is minutes, not
seconds.
\end{description}

We model these as a per-rule \emph{tier} (\texttt{stream} or \texttt{batch}) with
a configurable latency added to every alert timestamp. Reporting a detection
latency that excludes pipeline delay would flatter the architecture; making the
delay an explicit, swept parameter (\S\ref{sec:results}) makes the claim honest
and lets a reader substitute latencies measured in their own account.

\subsection{Normalized event schema}

Every layer emits the same record: timestamp, source, event type, actor, target,
outcome, source address, and a rule-specific attribute map. The schema is a small
subset of the common event models used in practice --- enough to write portable
detections, small enough that its data dictionary fits on a page.

Each event additionally carries a ground-truth label (\texttt{benign} or an attack
identifier). The label is written by the scenario, never read by a rule --- a
property enforced by a test that fails if the string appears anywhere in the rule
module --- and used only for scoring.

\subsection{Deployment}

The cloud instantiation is Terraform for a stream, an encrypted archive with
lifecycle transitions, an audit trail with object-level data events on the
archive, two functions, an event bus, findings aggregation, and least-privilege
roles in which the detector can observe but not act and the responder can act
only through a narrow set of reversible containment operations. The detection
function imports the same rule package the experiments run, shipped as a layer:
a single implementation is what allows results measured offline to be claimed for
the deployed system.

\subsection{Automated response, deliberately constrained}

Only containment is automated, and only with reversible actions. A ground segment
controls a vehicle that cannot be taken offline for maintenance; an automated
action that safes a spacecraft in response to a false positive causes a real
mission outage, and the response becomes the incident. The most aggressive action
the automation can take is to stop the \emph{ground segment} from transmitting.
Nothing in the response path commands the spacecraft.

\section{Detection content}
\label{sec:detection}

Eleven rules cover the five scenarios. They fall into three kinds, and the
distinction is not cosmetic --- the kinds behave differently under both attack and
noise.

\begin{description}
\item[Policy rules] fire on a violation of a control that nominal operations
never produce: a command refused by an allowlist, a frame refused by the
vehicle, an audit trail stopped. They are the cheapest and most precise rules in
the pack and, in our results, the ones that carry most of the detection load.
\item[Thresholded rules] aggregate an individually noisy signal over a window: a
burst of failed authentications, a cluster of frame integrity failures, a count
of distinct objects read. Their threshold trades latency against noise directly.
\item[Behavioural rules] compare against a profile learned during a warm-up
window: an unprofiled network origin, a principal issuing its first critical
command, an hourly read volume far above baseline. They are the only rules that
can catch an adversary who violates no policy.
\end{description}

\begin{table}[t]
\centering
\small
\caption{Detection content. Tier determines latency: \texttt{stream} rules see an
event within seconds, \texttt{batch} rules cannot beat their own aggregation
window.}
\label{tab:rules-catalogue}
\begin{tabular}{@{}llll@{}}
\toprule
ID & Kind & Tier & Detects \\
\midrule
R01 & thresholded  & stream & A1 \\
R02 & behavioural  & stream & A1, A2, A3 \\
R03 & policy       & stream & A2 \\
R04 & behavioural  & stream & A2 \\
R05 & policy       & stream & A2, A4 \\
R06 & thresholded  & stream & A3 \\
R07 & behavioural  & batch  & A3 \\
R08 & policy       & stream & A4 \\
R09 & thresholded  & batch  & A5 \\
R10 & statistical  & batch  & A5 \\
R11 & policy       & stream & A3, A5 \\
\bottomrule
\end{tabular}
\end{table}

Every rule maps to a SPARTA technique \cite{sparta} and, where an analogue
exists, to ATT\&CK \cite{attack}. No attack depends on a single rule; that is a
design requirement, and the ablation in \S\ref{sec:results} exists partly to
check that removing one control does not open a blind spot.

\subsection{Two rules worth describing in detail}

\paragraph{R06, frame integrity, and why it is thresholded.}
A radio link corrupts frames on its own, particularly at low elevation at the
edges of a pass. In our model the frame error probability rises exponentially as
elevation falls, producing roughly six to nine corrupted frames per week out of
about \num{1800}. A rule that alerts on a \emph{single} failed tag detects
tampering approximately \SI{60}{\second} sooner and generates enough noise to be
switched off within a week --- the exact failure mode a detection architecture
exists to avoid. R06 therefore requires three failures within
\SI{120}{\second}. \S\ref{sec:results} quantifies both sides of that trade.

\paragraph{R07, physics, and the rail exclusion.}
R07 checks three things about each telemetry channel: that it lies inside its
physical envelope, that it is not changing faster than the subsystem can change,
and that it is not frozen. The third check is what catches a masking attack that
holds a degrading measurement constant. Its first implementation was the largest
single source of false positives in the testbed, because a fully charged battery
in sunlight is legitimately constant. Excluding channels resting against a
physical rail removed that entire class. The lesson generalises beyond this
testbed: a stuck-value detector on any physical process must know which constants
are physics and which are faults.

\paragraph{Analytics should not run on bits already known to be corrupt.}
R07's second-largest false-positive source was subtler than the rail problem and
more instructive. A corrupted frame contains implausible values by definition, so
the physics rule dutifully alerted on link noise that the integrity check had
\emph{already} rejected --- two rules alarming on one event, one of them
uselessly. Excluding samples whose frame failed its tag check removed every such
alert and left a clean division of labour: R06 owns provenance failures, R07 owns
content. The general form of the lesson is that a detection pipeline should
establish trust in an input before reasoning about its meaning.

\subsection{Portability}

Each rule exists in two synchronised forms: executable Python, which both the
experiments and the deployed function import, and Sigma \cite{sigma}, which can
be converted for another backend. A test asserts that the two sets of rule
identifiers are equal, so the packs cannot silently drift apart --- a failure mode
common enough in practice to be worth automating against.

\section{Methodology}
\label{sec:methodology}

\subsection{Protocol}

One trial has two phases:

\begin{enumerate}
\item \textbf{Warm-up} (\ResWarmupDays{} simulated days). Nominal operations with
the engine in learning mode. Behavioural rules build their profiles and emit
nothing.
\item \textbf{Measurement} (\ResMeasureDays{} simulated days,
\ResEventsPerWindow{} events). Nominal operations continue while the five attacks
are injected at staggered start times, separated by \SI{20}{\hour} plus a
per-trial random jitter of up to \SI{8}{\hour}. The spacing makes every alert
attributable to at most one attack; the jitter prevents each trial from placing
attacks at the same point of the orbital cycle, which would report schedule
artefacts as variance.
\end{enumerate}

We run \ResTrials{} independent trials with different seeds and report means with
\SI{95}{\percent} normal-approximation confidence intervals.

\subsection{The benign workload}

A testbed with no benign anomalies reports a false-positive rate of zero and says
nothing useful. Ours generates four classes of benign noise, chosen because each
is a real analyst's daily experience: operators mistyping passwords (occasionally
enough times in a row to resemble a spray); operators connecting from mobile or
travel networks; the radio link corrupting frames at low elevation; and a
legitimate weekly reprocessing job that bulk-reads the archive in a pattern
superficially identical to exfiltration. Their \emph{rates} are chosen, not
measured from a real mission --- a limitation we return to in
\S\ref{sec:limitations}.

\subsection{Scoring definitions}

\begin{description}
\item[True positive.] An alert at least one of whose evidence events carries a
non-benign ground-truth label. Attribution therefore follows causality, not
temporal coincidence: an alert that merely fires during an attack window, caused
by unrelated benign activity, counts as a false positive.
\item[Time to detect (TTD).] The first attributed alert's actionable timestamp
minus the attack's first malicious action. It includes the pipeline latency of
the tier that produced the alert. For the replay scenario the reference point is
the first re-transmission, not the passive radio capture, because capture leaves
no observable in the ground segment.
\item[Time to contain (TTC).] The first automated containment action on the
attack's entity minus the same reference point.
\item[Impact.] Malicious actions that took effect --- commands executed, objects
read, forged frames accepted --- reported both in total and, separately, as
critical commands executed and megabytes exfiltrated, because ``one executed
thruster firing'' and ``one object read'' are not comparable harms.
\item[Impact after containment.] The subset of the above occurring after the
first containment action: what the response did \emph{not} prevent.
\end{description}

\subsection{Deduplication as part of the system under test}

Without a per-entity cooldown, a single incident produces an alert storm and the
false-positive rate becomes meaningless. The cooldown (\SI{300}{\second} by
default) is therefore an explicit, swept parameter rather than a hidden
implementation detail.

\subsection{Ablation and sensitivity}

The ablation repeats the whole campaign with one control disabled at a time
across seven preventive and responsive mechanisms, plus an all-disabled
configuration. The sensitivity analysis sweeps four tuning parameters --- the
frame-integrity threshold, the batch-tier latency, the archive enumeration
threshold and the deduplication cooldown --- one at a time with everything else
at default.

\subsection{Reproducibility}

Every run is seeded and deterministic; a test asserts that two trials with the
same seed produce identical latencies and false-positive counts. Fourteen
simulated days take about two seconds of wall time. Every number in this paper is
generated from the experiment JSON by a script that emits the LaTeX macros and
table bodies the document includes, so a stale figure in the text is not possible.

\section{Results}
\label{sec:results}

\subsection{RQ1 --- detection and latency}

All five attacks were detected in all \ResTrials{} trials
(Table~\ref{tab:baseline}). Mean time to detect across the campaign is
\ResTTDMean{}\,s.

\begin{table*}[t]
\centering
\footnotesize
\caption{Per-attack detection performance over \ResTrials{} trials. TTD and TTC
are measured from the first malicious action and include pipeline latency.
``Actions that succeeded'' counts adversary actions that took effect, whether or
not they were detected. Median and p95 are given alongside the mean because A2's
latency is bimodal, and for that attack the mean describes no trial that actually
occurred; its maximum is \ResTTDMaxAtwo{}\,s.}
\label{tab:baseline}
\begin{tabular}{@{}lrrrrrrlr@{}}
\toprule
Attack & Det.\ (\%) & Mean TTD (s) & 95\,\% CI & Median TTD (s) & p95 TTD (s) &
Mean TTC (s) & First rule & Actions succeeded \\
\midrule
\inputrows{generated/tab-baseline.tex}
\bottomrule
\end{tabular}
\end{table*}

The striking feature is not the latency but its \emph{variance}, and what
governs it. For four of the five attacks the distribution is tight across trials
with randomised start times: A4 is identical in every trial, and A1, A3 and A5
vary by $\pm$\ResCIAone{}\,s, $\pm$\ResCIAthree{}\,s and $\pm$\ResCIAfive{}\,s
about means of tens to hundreds of seconds. For those four, latency is a property
of the architecture rather than of the adversary, and it decomposes cleanly into
two terms.

A2 is the exception, and it is the more instructive case. Its mean of
\ResTTDAtwo{}\,s describes no trial that occurred, because the distribution is
bimodal. In \ResFirstRuleTrialsAtwo{} of \ResTrials{} trials the hijacked session
is observable, \ResFirstRuleAtwo{} fires on it, and the attack is detected in
\ResTTDMedianAtwo{}\,s. In the remaining \ResFallbackTrialsAtwo{} the session
never materialises: the operator account A2 targets was already blocked by an
earlier automated response --- usually one triggered by a false positive days
before --- and the API refuses the hijack without emitting an authentication
event at all. With no logon to reason about, \ResFirstRuleAtwo{} has nothing to
evaluate, and detection falls through to \ResFallbackRuleAtwo{} on the RF phase,
which can only fire while the spacecraft is in a pass window. Latency is then set
by orbital mechanics rather than by the pipeline, reaching \ResTTDMaxAtwo{}\,s.

Two things follow. First, a campaign of ten trials has a substantial chance of
drawing none of the slow mode and reporting a confidence interval of zero for A2,
which is why the baseline here runs \ResTrials{}. Second, the second mode is not
an adversary capability: it is an interaction between the response system and the
detection system, in which containing one attack degrades the observability of
the next. We return to this in Section~\ref{sec:discussion}.

The first term is the pipeline. A2 (in its fast mode), A3 and A4 are detected in
\ResTTDMedianAtwo{}\,s, \ResTTDAthree{}\,s and \ResTTDAfour{}\,s respectively ---
essentially the stream-path latency itself --- because their first malicious
action is already sufficient evidence: a session from an unprofiled origin, a
frame that failed its tag, a duplicated frame counter. Nothing needs to
accumulate.

The second term is the rule's own threshold, and it dominates whenever evidence
must accumulate. A1 is detected at \ResTTDAone{}\,s not because the pipeline is
slow but because R01 waits for five failed authentications, and the adversary
sprays one every eleven seconds. A5 is detected at \ResTTDAfive{}\,s because R09
waits for twenty-five distinct objects at one every four seconds. In both cases
the SOC could halve the latency by halving the threshold, and would pay for it in
noise --- which is precisely the trade the sensitivity analysis below quantifies.

\subsection{RQ2 --- false positives}

Across the measurement windows the pack produced
\ResFPDay{}~$\pm$~\ResFPDayCI{} false positives per day at \ResPrecision{}\,\%
precision (Table~\ref{tab:rules}). For a single spacecraft that is between three
and four alerts per week an analyst will close as benign. Whether this is
acceptable is not a question about precision in the abstract: it is a question
about whether a team that already runs passes can absorb one extra triage per
weekday, and at this rate it plainly can.

\begin{table}[t]
\centering
\footnotesize
\caption{Per-rule alert quality, summed over all trials. Policy rules produce no
false positives at all; every false positive in the pack comes from a
thresholded or behavioural rule.}
\label{tab:rules}
\begin{tabular}{@{}llrrrr@{}}
\toprule
ID & Rule & TP & FP & Prec. & FP/day \\
\midrule
\inputrows{generated/tab-rules.tex}
\bottomrule
\end{tabular}
\end{table}

The distribution matters more than the total. Policy rules --- a command outside
an allowlist, a frame the vehicle refused, an audit trail stopped --- produced no
false positives whatsoever, because in a correctly configured ground segment
nominal operations genuinely never generate those events. All false positives came
from the thresholded and behavioural rules, and they came from precisely the
benign anomalies the workload was built to contain: an operator whose typos
crossed the brute-force threshold, and the weekly reprocessing job, whose access
pattern R09 cannot distinguish from exfiltration on volume alone.

One rule never fired at all. R10, the archive volume anomaly, produced neither
true nor false positives: in benign traffic no principal exceeded four standard
deviations of its own hourly baseline, and during A5 the enumeration rule and the
automated containment both acted well before the first hourly bucket closed. We
report this rather than quietly dropping the rule. It is redundant \emph{in this
configuration}, and it would stop being redundant against an adversary who reads
a small number of very large objects --- exactly the case an object-count
threshold cannot see. A rule that never fires is either dead content or an
uncollected insurance premium, and distinguishing the two requires saying which
attack it is for.

R07's remaining false positives are the most interesting result in this table,
because the attack causes them. Once A3 disables frame authentication, every
subsequent frame arrives unauthenticated, and the ground segment loses its
ability to tell link corruption from injection. The physics rule --- now the only
detection left --- therefore inherits the radio's noise and starts alerting on
naturally corrupted frames. Disabling frame authentication does not merely remove
a detection; it degrades the precision of the detection that replaces it. We
report this rather than suppressing it, because a reader deploying content
analytics as a fallback for cryptographic integrity should know that the fallback
is noisier precisely when it is needed.

The reprocessing-job case deserves a direct statement rather than a footnote. The correct
remedy is to exclude the reprocessing principal by identity and alert instead when
\emph{that} principal reads from an unexpected origin. Raising the threshold
until the job stops alerting --- the tempting fix --- would also raise it past the
attack, which is how detections quietly stop working.

\subsection{RQ3 --- which logs identify each attack}

Table~\ref{tab:evidence} maps each attack to the log source that carried its
first evidence. No attack was visible in only one source, but the sources are not
interchangeable: three of the five were first identified by ground-segment logs
that no cloud-native service produces by default. A deployment that instruments
only the cloud control plane would have detected A1 and A5, would have detected
A3 only at its second phase, and would have missed A2's radio variant and A4
entirely.

\begin{table}[t]
\centering
\footnotesize
\caption{First evidence by log source. Sources marked \emph{ground} are produced
by the ground segment itself and must be shipped deliberately.}
\label{tab:evidence}
\begin{tabular}{@{}lll@{}}
\toprule
Attack & First evidence & Origin \\
\midrule
A1 & \texttt{api.auth} failures + success & ground \\
A2 & \texttt{api.auth} session origin, then \texttt{gs.uplink} & ground \\
A3 & \texttt{gs.downlink} integrity failures & ground \\
A4 & \texttt{gs.uplink} replay rejection & ground \\
A5 & \texttt{cloud.data} object reads & cloud \\
\bottomrule
\end{tabular}
\end{table}

\subsection{RQ4 --- which defence did the work}

Table~\ref{tab:ablation} repeats the campaign with one control removed at a time.
Detection rate is not the interesting column: every configuration still detects
every attack, which is the defence-in-depth property the pack was designed for.
Impact is the interesting column.

\begin{table*}[t]
\centering
\small
\caption{Defence ablation, \ResTrials{} trials per configuration. Detection
survives every single-control removal; impact does not.}
\label{tab:ablation}
\begin{tabular}{@{}lrrrrrr@{}}
\toprule
Configuration & Det.\ (\%) & TTD (s) & Actions & Critical cmds & MB out & FP/day \\
\midrule
\inputrows{generated/tab-ablation.tex}
\bottomrule
\end{tabular}
\end{table*}

Three results stand out.

\paragraph{Frame authentication and the anti-replay window are the only controls
that protect the vehicle.} With the full stack, \AblFullCrit{} critical commands
reach the spacecraft on average --- these are A3's \texttt{TLM\_AUTH\_DISABLE},
issued through a hijacked flight-director session and therefore inside the
allowlist, which is exactly the case an entitlement check cannot catch. Removing
frame authentication raises the figure to \AblNoSDLSCrit{}; removing the
anti-replay window raises it to \AblNoReplayCrit{}. Nothing in the cloud analytics tier substitutes for either.
Detection still fires --- promptly --- but the command has already executed on a
physical vehicle when it does. This is the clearest instance in our data of the
general principle that on a system with a physical actuator, detection is not a
substitute for prevention.

\paragraph{Automated containment prevents most of the impact.} Disabling it
raises total successful malicious actions from \AblFullSucc{} to
\AblNoResponseSucc{} and exfiltrated volume from \AblFullMB{}\,MB to
\AblNoResponseMB{}\,MB, with identical detection latency. The alerts were the
same; the difference is entirely in whether anything acted on them within
\SI{20}{\second}.

\paragraph{A control can be redundant because another one already covered it.}
Disabling MFA changed no outcome at all, which is initially surprising for the
one control that genuinely prevents A1. The explanation is in the timeline:
R01 fires on the fifth failed authentication and containment revokes the account
before the adversary's spray reaches the correct password. Detection substituted
for prevention here --- but only because the adversary was noisy. In the
all-disabled configuration, where nothing contains the account either, A1 becomes
the second most productive attack in the campaign. Redundancy between a
preventive and a detective control is worth having and worth knowing about; it is
not a reason to remove either.

\paragraph{Some controls cost nothing and are worth having anyway.} Removing the
rate limit or the pass-window feasibility check changed no outcome in our
scenarios. We report this rather than omitting it: their value in this testbed
is that they force an adversary into behaviour that is easier to detect, not that
they block anything on their own.

\subsection{Sensitivity of the tuning parameters}

Table~\ref{tab:sweep-integrity} shows the frame-integrity threshold trade-off
directly. At a threshold of one the rule detects tampering
\SweepIntegOneTTD{}\,s on average but produces \SweepIntegOneFP{} false
positives per day; at three it detects at \SweepIntegThreeTTD{}\,s with
\SweepIntegThreeFP{} per day. Buying roughly a minute of latency in exchange for
a rule an analyst will not mute is, in our judgement, the right trade --- but it
is a judgement, and the table is what makes it a judgement rather than a guess.

\begin{table}[t]
\centering
\footnotesize
\caption{Frame-integrity threshold: latency against noise.}
\label{tab:sweep-integrity}
\begin{tabular}{@{}rrrrr@{}}
\toprule
Threshold & Mean TTD (s) & Det.\ (\%) & FP/day & Prec.\ (\%) \\
\midrule
\inputrows{generated/tab-sweep-integrity-threshold.tex}
\bottomrule
\end{tabular}
\end{table}

The batch-tier latency sweep is starker still: moving from a
\SI{30}{\second} to a \SI{900}{\second} query cadence moves campaign mean TTD from
\SweepBatchFastTTD{}\,s to \SweepBatchSlowTTD{}\,s while
changing the false-positive rate not at all. Query cadence is free latency ---
the most consequential and most easily overlooked parameter in the architecture.

\begin{table}[t]
\centering
\footnotesize
\caption{Batch-tier query cadence. Latency is bought with money, not with noise.}
\label{tab:sweep-batch}
\begin{tabular}{@{}rrrrr@{}}
\toprule
Latency (s) & Mean TTD (s) & Det.\ (\%) & FP/day & Prec.\ (\%) \\
\midrule
\inputrows{generated/tab-sweep-batch-latency-s.tex}
\bottomrule
\end{tabular}
\end{table}

\subsection{RQ5 --- cost}

Table~\ref{tab:cost} gives the modelled monthly bill for the measured workload:
approximately USD~\CostMonthly{} for one spacecraft and one ground station, of
which \CostTopShare{}\,\% is \CostTopService{}.

\begin{table}[t]
\centering
\footnotesize
\caption{Modelled monthly cost, single spacecraft, on-demand list prices, no
committed-use discounts. Lines below USD~0.01 omitted.}
\label{tab:cost}
\begin{tabular}{@{}llr@{}}
\toprule
Service & Cost driver & USD/month \\
\midrule
\inputrows{generated/tab-cost.tex}
\bottomrule
\end{tabular}
\end{table}

\begin{table}[t]
\centering
\footnotesize
\caption{Cost scaling. The architecture is a fixed platform cost at small scale.}
\label{tab:cost-scaling}
\begin{tabular}{@{}rrrr@{}}
\toprule
Spacecraft & Events/day & USD/month & USD/spacecraft \\
\midrule
\inputrows{generated/tab-cost-scaling.tex}
\bottomrule
\end{tabular}
\end{table}

The scaling behaviour (Table~\ref{tab:cost-scaling}) is the practically important
result: per-spacecraft cost falls to USD~\CostPerHundred{} at a hundred vehicles.
At lab scale the bill is dominated by the minimum billable capacity of the search
tier rather than by event volume --- the architecture is not expensive because it
processes much, but because it must exist at all. For a single-satellite academic
or startup mission this is the dominant economic fact, and it argues for pooling
the analytics tier across missions rather than for reducing telemetry.

\section{Discussion}
\label{sec:discussion}

\subsection{Answering the research question}

Cloud-native security telemetry does detect all five attack classes, quickly and
at a false-positive rate a small team can absorb --- but the qualifier matters
more than the headline. Three of the five attacks were first identified by logs
the \emph{ground segment} produces, not by anything a cloud provider generates on
its own. The cloud tier supplied the pipeline, the correlation, the findings
store and the response automation; it did not supply the evidence. An
organisation that adopts cloud-native security tooling and instruments only what
that tooling emits by default will detect credential abuse and data
exfiltration, and will be blind to the attacks that reach the vehicle.

The honest form of the answer is therefore conditional: cloud-native security
telemetry is effective for ground segments \emph{if the ground segment is
instrumented as a first-class log source}. That conclusion is unglamorous and, in
our view, the most useful thing this study produces.

\subsection{Detection is not a substitute for prevention}

The ablation makes this concrete in a way that a narrative argument cannot.
Removing frame authentication does not reduce the detection rate at all --- every
forged command is still detected, promptly --- yet it raises the number of
critical commands executing on the spacecraft from \AblFullCrit{} to
\AblNoSDLSCrit{}. A perfect detector and a \SI{4}{\second} pipeline do not undo a
thruster firing.

This is a general property of systems with physical actuators, and it argues for
a specific division of labour. Preventive controls belong as close to the
actuator as possible --- on the vehicle, where the anti-replay window lives, not
in the analytics tier. Detection exists to catch what prevention does not cover
and to bound the duration of an intrusion, not to replace it. In a ground segment
the distinction is unusually sharp, because the window in which prevention can
act is the pass, and the pass is over in seven minutes.

\subsection{Cryptographic integrity and content analytics are complementary}

Attack A3 was designed to separate two defences that are often treated as
alternatives. Its first phase injects implausible values with a forged tag: the
cryptographic check catches it immediately. Its second phase first disables frame
authentication and then freezes a degrading measurement at a plausible constant.
After that point the tag check has nothing to say --- every frame is
``correctly'' unauthenticated --- and the only detection left is the one that
reasons about whether the measurement makes physical sense.

The measurement adds a wrinkle we did not anticipate. Once authentication is
off, the physics rule also inherits the link's own noise, because the ground
segment can no longer separate a corrupted frame from an injected one. The
fallback detection is therefore least precise exactly when it becomes the only
one --- a property worth knowing before treating content analytics as a
substitute for provenance.

Neither defence subsumes the other. Cryptography establishes provenance and says
nothing about content; physics analytics evaluates content and says nothing about
provenance. A ground segment that has only one of them has a documented blind
spot, and A3 traverses it in two moves.

\subsection{Where the false positives came from}

Every false positive in the campaign came from a thresholded or behavioural rule,
and each traced to a benign anomaly that a real ground segment produces every
week. This has a practical implication for anyone building such a pack: the
policy rules --- refused commands, refused frames, stopped audit trails --- are
close to free, and should be written first. They are precise because the events
they match genuinely do not occur in nominal operations, and that property comes
from the ground segment's small, well-defined operational envelope rather than
from clever rule authoring.

The behavioural rules earn their noise, though. R02, the unprofiled-origin rule,
was the first detector for two of the five attacks. In a general enterprise,
origin profiling is a notoriously noisy control; in a ground segment, where the
operator population is a handful of named people working from a handful of known
networks, it is one of the strongest signals available. Threat detection content
does not transfer between domains unchanged, and this is a case where the domain
makes a weak enterprise rule into a strong one.

\subsection{Containment degrades the observability of later attacks}

The A2 latency distribution in Section~\ref{sec:results} is bimodal for a reason
that has nothing to do with the adversary. Automated containment blocks a
principal, and in this architecture a block is permanent for the remainder of the
campaign. When a false positive blocks a benign operator --- which happens at the
false-positive rate the pack already reports --- that account stops being usable,
and the mission-control API refuses any later session for it \emph{silently},
without emitting an authentication event. An adversary who subsequently targets
that account therefore produces no logon for the behavioural rules to evaluate.
Detection does not fail, but it falls back to a slower path: in our campaign,
from \ResTTDMedianAtwo{}\,s on the session to as much as \ResTTDMaxAtwo{}\,s on
the radio-frequency phase, which can only be observed during a pass.

This is worth stating plainly because it inverts an assumption that is easy to
make. Automated response is normally evaluated on what it prevents, and
Section~\ref{sec:results} reports that in those terms. But a response action also
changes what the telemetry contains, and a control that removes an actor from the
system removes that actor's events along with it. The effect is not confined to
this testbed: any architecture that disables accounts automatically and logs
refusals at a different fidelity than successes will have the same blind spot.
Two mitigations follow directly --- expire automated blocks rather than leaving
them indefinite, and emit an explicit event when an action is refused because the
principal is blocked, so that the refusal is itself a detectable signal. We did
not implement either here; measuring them is future work.

\subsection{Implications for practitioners}

\begin{enumerate}
\item \textbf{Ship ground-segment logs.} Authentication, telecommand submission,
uplink verdicts and frame integrity are the evidence for three of five attacks.
Nothing generates them for you.
\item \textbf{Put the anti-replay window on the vehicle and verify it.} It is the
only control in the study that both prevents an attack and produces its alert.
\item \textbf{Choose your query cadence deliberately.} Moving batch-tier latency
from \SI{900}{\second} to \SI{30}{\second} improved campaign mean latency by more
than any rule change we made, at no cost in false positives.
\item \textbf{Exclude known-good service principals by identity, never by
threshold.} The weekly reprocessing job and the exfiltration attack are
indistinguishable on volume; raising the threshold past the job also raises it
past the attack.
\item \textbf{Automate containment, not remediation.} Containment cut impact by
roughly four fifths in our data. Automating anything that commands the vehicle
would make a false positive into a mission outage.
\item \textbf{Budget for the platform, not the volume.} At single-mission scale
the cost is a floor, not a function of telemetry. Pool the analytics tier across
missions rather than reducing what you collect.
\end{enumerate}

\section{Limitations and threats to validity}
\label{sec:limitations}

\subsection{Construct validity}

The most significant limitation is that the offline harness \emph{models}
pipeline latency rather than measuring it in a live cloud account. Our
stream-path figure of \SI{4}{\second} and batch-path figure of \SI{150}{\second}
are plausible for the services in question but they are parameters, not
observations. We mitigate this in three ways: the parameters are named rather
than buried; the sensitivity analysis reports results across a wide range of
them; and the deployed functions import the identical rule package, so a reader
with an account can substitute measured latencies without changing the detection
logic. Absolute latencies should nonetheless be read as \emph{lower bounds
conditioned on the modelled pipeline}, and the relative comparisons --- stream
against batch, threshold against threshold --- are the durable results.

The impact measure is coarse. Counting ``actions that succeeded'' treats an
executed thruster firing and a read object as one unit each, which is why we
report critical commands and exfiltrated volume separately. A severity-weighted
composite would be more informative and would also be more arbitrary; we chose
transparency.

\subsection{Internal validity}

Ground-truth labels are written by the attack scenarios themselves, so an error
in a scenario would produce a systematically wrong score. We guard against the
most dangerous form of this --- a rule reading the label it is being scored
against --- with a test that fails if the label appears anywhere in the rule
module. Attribution by evidence causality rather than by time window removes the
other common failure, in which unrelated benign alerts during an attack window
are counted as detections.

Attacks are separated by at least \SI{20}{\hour} so that alerts attribute
cleanly. Real intrusions overlap, and a campaign of simultaneous attacks would
stress deduplication and analyst triage in ways this design does not measure.

\subsection{External validity}

The benign workload is synthetic. Its noise sources are realistic in kind ---
typos, roaming operators, link fades, a weekly bulk job --- but their rates are
chosen. The false-positive figures should therefore be read as
\emph{order-of-magnitude} results conditioned on that workload, not as a
prediction for a specific mission. A team adopting this pack should expect to
re-tune thresholds against their own traffic; the sweep tables show how much
movement each threshold buys.

The configuration is one spacecraft, one ground station, one cloud account.
Multi-tenant ground-station-as-a-service, cross-mission analytics, and
constellation-scale operations each introduce isolation and correlation questions
this study does not address.

The spacecraft model is deliberately shallow: circular two-body propagation, no
perturbations, no manoeuvre dynamics, no link budget in decibels. This is
sufficient because the detections depend on pass timing, physical envelopes and
frame outcomes, all of which the model reproduces --- but it means the testbed
cannot be used to study attacks whose detection depends on finer physics, such as
subtle orbit-determination manipulation.

\subsection{What we did not evaluate}

Availability attacks (jamming, uplink flooding) are excluded because they are
detected by link metrics rather than security telemetry. Supply-chain compromise
of flight software, physical attacks on the ground station, and adversaries who
obtain the SDLS key are excluded by the threat model. Analyst experience --- the
question of whether \ResFPDay{} false positives per day is actually tolerable in
a given SOC --- is a human-factors question our method cannot answer.

\section{Conclusion}

We built a reproducible testbed spanning the full satellite ground segment,
executed five controlled attacks against it, and measured what a cloud-native
security architecture actually sees. All five attacks were detected across
\ResTrials{} trials, with a campaign mean time to detect of \ResTTDMean{}\,s,
\ResFPDay{}~$\pm$~\ResFPDayCI{} false positives per day, \ResPrecision{}\,\%
precision, and a modelled cost of approximately USD~\CostMonthly{} per month for
a single spacecraft, falling to USD~\CostPerHundred{} per spacecraft at a hundred.

The answer to the research question is a qualified yes. Cloud-native security
telemetry detects these attacks well, but the cloud tier contributes the
pipeline, the correlation and the response --- not the evidence. Three of the
five attacks are visible only in logs the ground segment must be instrumented to
produce. Detection latency is set by which tier a rule runs in rather than by
adversary behaviour, so architectural placement dominates rule tuning.
Cryptographic frame integrity and physics-aware analytics defend against
different halves of a tampering attack and neither substitutes for the other. And
on a system whose actuator is a spacecraft, detection does not substitute for
prevention: removing frame authentication changed our detection rate not at all
while allowing critical commands to execute on the vehicle.

\subsection*{Future work}

Four directions follow directly. First, replace the modelled pipeline latencies
with measurements from a deployed account, which the shared rule package makes a
matter of instrumentation rather than reimplementation. Second, replace the
synthetic benign workload with traces from an operating mission --- the single
change that would most improve the external validity of the false-positive
results. Third, extend to constellation scale and to multi-tenant
ground-station-as-a-service, where cross-mission correlation becomes possible and
tenant isolation becomes a new attack surface. Fourth, study overlapping and
adaptive adversaries, including one that observes the containment response and
adjusts --- a case the current staggered design deliberately excludes.

\subsection*{Availability}

The testbed, the eleven detection rules in both executable and Sigma form, the
five attack scenarios, the incident-response playbooks, the Terraform for the
cloud deployment, the labelled event dataset and the analysis scripts that
generate every number in this paper are released at
\url{https://github.com/rikosoo/space-ground-station-security-lab}. All results
in this paper are regenerated by a single command and are deterministic under a
fixed seed.


\balance
\bibliographystyle{plain}
\bibliography{refs}

\end{document}
